% Physical Review paper (PRL, PRB, PRX, PR Applied, PR Materials and other APS journals).
% Uses the revtex4-2 class that TeX Live ships; the options choose the journal and layout:
% prl / prb / prx / prapplied / prmaterials, and preprint (one column) or reprint (as printed).
% See APS's author guide (apsguide4-2.pdf on CTAN) for the full list.
\documentclass[aps,prb,reprint,superscriptaddress,amsmath,amssymb]{revtex4-2}

\usepackage{graphicx}
\usepackage{booktabs}
\usepackage[colorlinks=true,allcolors=blue]{hyperref}

\begin{document}

\title{Ion transport in a concentrated electrolyte near a lithium-metal interface}

\author{First Author}
\email{first.author@example.edu}
\affiliation{Department of Physics, Example University, City 100000, Country}
\author{Second Author}
\affiliation{Department of Physics, Example University, City 100000, Country}
\affiliation{Institute of Materials, Example Laboratory, City 100000, Country}

\date{\today}

\begin{abstract}
Abstracts in Physical Review journals are a single paragraph of up to 600 characters for
PRL and 500 words elsewhere. State the problem, the approach, and the main quantitative
result. This sample shows the parts of a paper: sections, an equation, a table and citations.
\end{abstract}

\maketitle

\section{Introduction}
Concentrated electrolytes change the solvation of lithium ions and the chemistry of the solid
electrolyte interphase~\cite{yamada2019,xu2004}. Here we model transport through such an
interphase with a one-dimensional drift--diffusion equation.

\section{Model}
The lithium concentration $c(x,t)$ obeys
\begin{equation}
  \frac{\partial c}{\partial t} = D \frac{\partial^2 c}{\partial x^2}
  - \frac{D e}{k_B T} \frac{\partial}{\partial x}\left(c \, E\right),
  \label{eq:transport}
\end{equation}
where $D$ is the diffusion coefficient and $E$ the electric field. Equation~(\ref{eq:transport})
reduces to Fick's law when $E = 0$.

\section{Results}
Table~\ref{tab:conductivity} lists the conductivity for three salt concentrations.

\begin{table}[b]
  \caption{Ionic conductivity $\sigma$ at 25\,$^\circ$C.}
  \label{tab:conductivity}
  \begin{ruledtabular}
    \begin{tabular}{lcc}
      Concentration (mol/kg) & $\sigma$ (mS/cm) & Transference number \\
      \colrule
      1.0 & 9.8 & 0.35 \\
      3.0 & 6.1 & 0.48 \\
      5.5 & 2.4 & 0.62 \\
    \end{tabular}
  \end{ruledtabular}
\end{table}

\begin{acknowledgments}
We thank the members of the laboratory for discussions.
\end{acknowledgments}

\bibliography{refs}

\end{document}
